\section{引言与主要结果}
\label{sec:1}
能量与动量是物理学中的基本概念。然而在格点场论（量子场论发展最成熟的非微扰表述）中，
与之对应的诺特流，即能动量
张量~\cite{Callan:1970ze,Coleman:1970je,Freedman:1974gs,Joglekar:1975jm}，的构造却
并非直截了当~\cite{Caracciolo:1988hc,Caracciolo:1989pt}，因为平移不变性被格点结构
显式破缺。在近期的一篇论文~\cite{Suzuki:2013gza}中，人们基于 Yang--Mills 梯度流
（在格点规范理论的语境中也称 Wilson 流）~\cite{Luscher:2010iy,Luscher:2011bx,Luscher:2013cpa}
提出了一种可能的方法来规避格点上能动量张量所固有的这一困难。%
\footnote{在文献~\cite{Giusti:2010bb,Giusti:2012yj,Robaina:2013zmb,%
Giusti:2013sqa,Giusti:2013mxa,Giusti:2014ila}中，一种由移位边界条件出发定义格点
能动量张量的有趣方法得到了发展。文献~\cite{Suzuki:2012wx}中则提出了基于
$\mathcal{N}=1$ 超对称的一种方法。}关于梯度流的近期综述可参见
文献~\cite{Luscher:2013vga}；其在格点规范理论中的应用可参见
文献~\cite{Borsanyi:2012zs,Borsanyi:2012zr,Fodor:2012td,Fritzsch:2013je,%
Monahan:2013lwa,Shindler:2013bia,Bonati:2014tqa}。

文献~\cite{Suzuki:2013gza}的基本思想如下：\footnote{以下推理受到了 E.~Itou
和~M.~Kitazawa 所作开创性尝试（未发表）的启发。}考虑纯 Yang--Mills 理论。梯度流按照带
流时间~$t$ 的流方程［见下文式~\eqref{eq:(3.1)}］使裸规范场~$A_\mu(x)$ 发生变形，
这使得规范场位形在正流时间下变得“光滑”。可以证明，在微扰论的任意阶上，对任意严格为正
的流时间~$t$，流动规范场~$B_\mu(t,x)$ 的任何局域乘积在用重正化参数表达时都是紫外（UV）
有限的~\cite{Luscher:2011bx}。特别地，无需任何相乘重正化因子即可使这些局域乘积有限。
换言之，它们是重正化了的复合算符。这样的 UV 有限量应当是“普适的”，也就是说，在取去
调节子的极限下，它们不依赖于所选定的 UV 正规化方式。这暗示了这样一种可能性：以梯度流
为中间工具，可以在维数正规化——平移不变性在其中是明显的%脚注：当然，维数正规化的缺点在于它只在微扰论中有定义。
——下定义的复合算符，与可以用格点正规化进行非微扰计算的复合算符之间架起桥梁。%
\footnote{因此我们隐含地假定：流动场的有限性虽然只能在微扰论中被严格证明，
但在非微扰层面仍然成立。}遵循这一想法，文献~\cite{Suzuki:2013gza}推导出了一个公式，
把正确归一且守恒的能动量张量表示为由流动规范场构成的某个组合在 $t\to0$ 极限下的结果。
该公式提供了一种用格点正规化计算能动量张量关联函数的可能途径，因为公式中的普适组合应当
与正规化无关。一个有意思之处在于，公式中（普适）系数的小流时间行为可以借助渐近自由由
微扰论确定。这意味着只要格距足够细，就无须再对这些系数作非微扰确定。（实际上，这些系数
的非微扰确定可能相当有用；至于如何实现这种确定，已在文献~\cite{DelDebbio:2013zaa}中作过
研究。）尽管文献~\cite{Suzuki:2013gza}中公式的有效性——尤其是守恒律在连续极限下的
恢复——仍待仔细考察，但基于该公式对有限温度 Yang--Mills 理论的相互作用测度（迹反常）与熵
密度的测量~\cite{Asakawa:2013laa}已显示出令人鼓舞的结果；这一方法即便从实用角度看也似乎
是有前途的。

在文献~\cite{Suzuki:2013gza}中，该方法只针对纯 Yang--Mills 理论建立。于是自然会问：
该方法能否推广到包含物质场、特别是费米子场的规范理论，从而获得更广的应用。本文正是要完成
这一推广。我们考虑一个矢量型规范理论\footnote{我们假定该理论具有渐近自由。}，其规范群~$G$
包含
\begin{equation}
   \text{$N_{\mathrm{f}}$ 个处于规范表示 $R$ 的 Dirac 费米子}.
\label{eq:(1.1)}
\end{equation}
为简单起见，我们假定所有 $N_{\mathrm{f}}$ 个费米子具有共同的质量；这一限制可以适当地放宽。

我们的主要结果是式~\eqref{eq:(4.70)}，后续各节将逐步给出该主公式的推导。对于主要只关心
最终结果的读者，这里简要说明如何阅读主公式~\eqref{eq:(4.70)}：左边是正确归一且守恒的能动量
张量（已减去真空期望值）；我们的公式~\eqref{eq:(4.70)}只在能动量张量与其他算符在位置空间
关联函数中相互分离的情形下成立。右边的组合由式~\eqref{eq:(4.1)}--\eqref{eq:(4.5)}
定义。其中 $G_{\mu\nu}^a(t,x)$ 与~$D_\mu$ 分别是流动规范场的场强与协变导数（流动规范场的
定义与
文献~\cite{Luscher:2010iy,Luscher:2011bx,Luscher:2013cpa}中的相同）；另一方面，我们的
\emph{加圈\/}流动费米子场
$\mathring{{\chi}}(t,x)$ 与~$\mathring{\Bar{\chi}}(t,x)$ 与
文献~\cite{Luscher:2013cpa}中的 $\chi(t,x)$
和~$\Bar{\chi}(t,x)$ 略有不同，二者通过式~\eqref{eq:(3.20)}
和~\eqref{eq:(3.21)} 相联系。式~\eqref{eq:(4.70)} 中的系数函数 $c_i(t)$ 由式%
~\eqref{eq:(4.72)}--\eqref{eq:(4.76)} 给出。那里的 $\Bar{g}(q)$ 与
$\Bar{m}(q)$ 分别是由
式~\eqref{eq:(4.20)} 和~\eqref{eq:(4.21)} 定义的跑动耦合与跑动质量参数；在全文中，$m$
表示费米子的（共同）重正化质量。
式~\eqref{eq:(4.72)}--\eqref{eq:(4.76)} 对应最小减除
（$\text{MS}$）方案，修正最小减除
（$\overline{\text{MS}}$）方案的结果可通过替换~\eqref{eq:(4.77)} 得到。$b_0$
和~$d_0$ 分别是重正化群函数，即式~\eqref{eq:(2.16)}
和~式~\eqref{eq:(2.18)} 的首项系数。能动量张量由
式~\eqref{eq:(4.70)} 右边的 $t\to0$ 极限给出。如上文纯 Yang--Mills 情形所述，右边的
组合是 UV 有限的，因而可以用格点正规化来计算其关联函数。这样便得到正确归一且守恒的
能动量张量的关联函数。然而，为了保证“普适性”，必须先取连续极限、再取 $t\to0$
极限。实际操作中，由于格距~$a$ 有限，流时间~$t$ 不能任意地小，以便保持与连续物理的联系；
于是存在自然约束
\begin{equation}
   a\ll\sqrt{8t}\ll R,
\label{eq:(1.2)}
\end{equation}
其中~$R$ 表示典型的物理尺度，例如强子尺度或盒子尺寸。因此 $t\to0$
的外推一般要求足够细的格点。

下面是我们对二次 Casimir 算符的定义：我们把表示~$R$ 的反厄米生成元~$T^a$
归一化为~$\tr_R(T^aT^b)=-T(R)\delta^{ab}$ 以及~$T^aT^a=-C_2(R)1$。同时记
$\tr_R(1)=\dim(R)$。由 $[T^a,T^b]=f^{abc}T^c$ 中的结构常数定义
$f^{acd}f^{bcd}=C_2(G)\delta^{ab}$。例如，对 $SU(N)$ 的基本
表示~$N$（此时 $\dim(N)=N$），常规归一化为
\begin{equation}
   C_2(SU(N))=N,\qquad T(N)=\frac{1}{2},\qquad
   C_2(N)=\frac{N^2-1}{2N}.
\label{eq:(1.3)}
\end{equation}
